\section{What we thought vs.\ what we measured}
\label{sec:honest}

A consolidated honesty pass on this engine's GPU story. Each
finding is paired with the prior we held and the measurement that
revised it.

\subsection{``GPU SVD wins because the reconstruction kernel is fast''}

\paragraph{Prior:} The SVD reconstruction microbench measures
$8$--$13\,\times$ per-XS-lookup speedup over a single-threaded CPU
loop. We expected this to translate to a $\sim$5--10$\,\times$
end-to-end GPU SVD speedup over CPU SVD for realistic transport.

\paragraph{Measurement (\texttt{outputs/method\_comparison\_2026-05-08.txt}):}
On Godiva, GPU SVD vs.\ 20-core rayon CPU SVD on a workstation Intel
mobile: $0.77\,\times$. The GPU SVD path is $1.3\,\times$
\emph{slower} than the parallel CPU path on integrated transport.
On the same problem against GPU pointwise table (apples to apples
on the same backend): also $0.77\,\times$. The launch overhead +
divergent memory access dominates the per-XS-lookup speedup at the
problem scale where Godiva sits.

\paragraph{Lesson:} Microbench speedups do not survive integration
when (a) the kernel does $\mathcal{O}(\text{microbench unit})$
arithmetic per the much-larger geometry-traversal + collision-
sampling work envelope, and (b) the per-XS-lookup pattern is
divergent within a warp.

\subsection{``Recursive transport is 6.74$\,\times$ on GPU''}

\paragraph{Prior:} Documented in earlier session notes, propagated
into prior STATUS.md headlines, cited in conversations as a
general ``GPU is faster than CPU'' fact.

\paragraph{Measurement:} The $6.74\,\times$ number is real but
scope-tagged. It comes from
\texttt{const\_xs\_transport\_persistent} --- a constant-XS
microbench that exercises the recursive geometry walker plus
fission banking via \texttt{atomicAdd}, but skips all the SVD /
table XS evaluation, all the angular distribution sampling, and
the energy-distribution sampling. As a measurement of how fast
\emph{geometry traversal} runs on the GPU, it is honest. As a
proxy for ``how fast our full transport runs on GPU'' it is
misleading.

\paragraph{Lesson:} Always carry the scope tag.
\texttt{[micro]}, \texttt{[godiva]}, \texttt{[assembly]} are not
decoration; they signal whether the number generalises. The
audit in this paper's preparation revised every header table in
\texttt{STATUS.md} and \texttt{README.md} to make this distinction
explicit.

\subsection{``Refill is 2$\,\times$ ICSBEP-wide''}

\paragraph{Prior:} After the first refill measurement at HMF-008
at 250\,k particles showed a clean $2\,\times$ throughput gain, we
expected the same uplift across the sweep.

\paragraph{Measurement (§\ref{sec:refill}):} The refill gain is a
function of where on the saturation curve the workload sits. At
5\,000 particles on RTX A1000, throughput is $0.987\,\times$ ---
refill adds overhead. At 250\,k on RTX 3080 (mid-curve),
$2.0\,\times$. At $\geq$1\,M on RTX 3080 (saturated), the
gain shrinks to single-digit percent. There is no universal
refill factor; the auto-recommender that reads device attributes
and routes by regime is the necessary engineering, not the
refill kernel itself.

\paragraph{Lesson:} Optimisations that exploit a specific
occupancy regime need device-aware dispatch. Static defaults are
either wrong or under-utilised on most cards.

\subsection{``CAB H-in-H2O fixes the HEU-COMP-INTER benchmark family''}

This finding is from the ENDF/B-VII.1 vs.\ VIII.1 transition that
happened during this paper's preparation and is documented in
\texttt{STATUS.md} and a recent commit. We include it here because
the diagnostic process used the same observability-first approach.

\paragraph{Prior:} VIII.1 ships the CAB H-in-H$_2$O thermal
scattering re-evaluation. We expected this to pull the
HEU-COMP-INTER-003 family closer to the handbook k=1.000 because
``better thermal scattering data should improve thermal-spectrum
benchmarks''.

\paragraph{Measurement (7-case A/B, CPU runner, 5 seeds, 150
batches, 40 inactive, 100\,k particles on RTX A1000 host):} VIII.1
shifts $k$ \emph{upward} by $+820 \pm 153\,$pcm uniformly across
the family. Cases that were marginally PASS under VII.1 became
FAIL under VIII.1 (case-2, case-3 against the handbook); cases
that were below 1.0 under VII.1 crossed zero into more
comfortable PASS (case-5, -6, -7).

\paragraph{Root cause via h5py probe:} VIII.1's CIELO-era U-235
re-evaluation raised $\sigma_f$ at the intermediate-spectrum
peak (100\,eV) by $+5.3\,\%$ and dropped $\sigma_\gamma$ by
$-5.5\,\%$, so the capture/fission ratio $\alpha$ dropped by
$-10.2\,\%$. The benchmark family's flux peaks at exactly that
energy. The CAB H-in-poly TSL change was not the dominant driver
here --- the H-in-CH$_2$ TSL was not even bound on the loaded
geometry (the only TSL binding in this family is
\texttt{c\_Be}).

\paragraph{Lesson:} When a benchmark family shifts under a
library change, the right diagnostic is to probe the library
directly (h5py at known-resonance energies) before assigning
the shift to a specific physics mechanism. The intuitive
hypothesis (``the TSL update fixed it'') was specifically wrong
about which TSL even applied.

\subsection{``Cache machinery just needs the right key''}

\paragraph{Prior:} Implementing a cache is straightforward; the
key design is the only interesting bit.

\paragraph{Measurement:} The cache itself is straightforward, but
\emph{verifying} the cache fires at runtime turned out to be
the harder problem. The engine's stdout was captured by the
Python sweep driver and never reached the log file. Without
\texttt{-{}-debug-metrics}, an operator running an ICSBEP sweep
has no way to falsify ``the cache is working''. We spent multiple
sessions reading code and convincing ourselves the cache must be
firing; the actual confirmation needed a $\sim$200-line
observability path.

\paragraph{Lesson:} Design the falsification path at the same
time as the optimisation. Cache hits are an absence of work, and
absences of work are hard to see without explicit telemetry.
